\documentclass[11pt,a4paper]{article}
\usepackage[T1]{fontenc}
\usepackage[utf8]{inputenc}
\usepackage{amsmath,amssymb,xcolor}
\usepackage[margin=25mm]{geometry}
\usepackage[colorlinks=true,urlcolor=blue]{hyperref}
\definecolor{nexorablue}{HTML}{2459B5}
\setlength{\parindent}{0pt}
\setlength{\parskip}{7pt}
\setlength{\emergencystretch}{2em}
\renewcommand{\familydefault}{\sfdefault}
\begin{document}

{\small\bfseries\color{nexorablue} NEXORA CPGE}\par
{\small MP / NEX-TOP-03 / Énoncé}\par
{\LARGE\bfseries Au-dessus d’une parabole tronquée}\par
\vspace{0.25em}\hrule\vspace{0.6em}

Dans $\mathbb R^2$, soit $A=\{(x,y):0<x\le2,\ y>x^2\}$. Déterminer l’intérieur, l’adhérence et la frontière de $A$. Donner un exemple de point de frontière appartenant à $A$ et un autre qui lui est extérieur.


\end{document}
